Un long solénoïde de longueur \qty{\v{L}}{cm} et de rayon \qty{\v{r}}{cm} compte $N = \v{N}$ spires. Il est parcouru par un courant de \qty{\v{I}}{A}. ($\mu_0 = 4\pi \times \qty{e-7}{T.m/A}$)

\begin{enumerate}[label=\alph*)]
    \item Déterminez le nombre de spires par mètre.
    \item Déterminez le champ magnétique à l'intérieur du solénoïde, loin des extrémités.
    \item Déterminez le flux magnétique à travers une seule spire.
    \item Quel courant faudrait-il faire circuler pour obtenir un champ de \qty{\v{Bt}}{mT} à l'intérieur ?
\end{enumerate}
